% 第9章 结语
\section{结语}\label{sec:conclusion}

把全文收成三句话，再指出三个未完成的方向。

\subsection{三个结论}

\textbf{第一，缝隙的存在性相对于观测装置而定义。}“这条信号有没有缝隙”是一个不完整的问句；完整的版本必须带上尺度与参照系。这不是修辞——第 \ref{sec:measurement} 章把它做成了严格的窗口条件：太密的网格被统计分辨力挡住（$\nu_0/\sqrt{\Delta}$），太疏的网格被信息衰减放走（$e^{-\lambda\Delta/2}$），两者之间才是可以动手的地带；理想装置极限下，窗口条件退化为原始的 $\delta_0>c$。窗口内部还有第二层结构：净收益的形状由毛收益的尺度衰减指数 $p$ 与成本的摊薄指数 $q$ 的赛跑决定，$q>p$ 才有内点最优，峰位对成本水平有 $1/(q-p)$ 的弹性。实证里，八年 A 股日线上的尺度扫描与这个结构大体相容：最慢的信号（非流动性）窗口宽而单调，最快的信号（1 日反转）被成本整体压制，中间带出现清晰的内点峰值（量比在 $\Delta=5$ 日、距均线与 60 日动量在 $\Delta\le3$ 日）。

\textbf{第二，缝隙的大小相对于参照系而定义。}$\alpha$ 的基准依赖性可以写成一行代数：$\alpha_{\mathcal{F}'}=\alpha_{\mathcal{F}}-\gamma\,\alpha_{\mathcal{F}}(g)$，消长量是“策略对新增工具的暴露”乘“新增工具自身的溢价”。完备参照系下缝隙归零——可交易缝隙的存在空间，恰是“不可被可交易工具复制的收益”这块地基。实践含义被压成一句诊断：\textbf{参照系扩张时缓慢衰减的是缝隙，骤然归零的是暴露。}

\textbf{第三，把存在性拆成四道可分别核验的门。}尺度、成本与实施、参照系、统计校正。每一道都能用真实数据回答，每一道的否定都有名有姓。四道全过才给“可交易”三个字；过不了的，卡在哪一道就写在哪一道——这份记录比结论本身更接近研究的增量。

\subsection{方法论上的三句实话}

其一，装置写不清，“有没有缝隙”就是伪问题。数据频率、执行延迟、成本参数、参照工具，这些不是附录里的实现细节，而是判据的自变量。

其二，核验要按“先便宜后昂贵、先否定后肯定”排序。估 $\hat\delta_0$ 与 $\hat c$ 的置信下界最便宜，尺度扫描次之，全链条的参照系分解与统计校正最贵。把便宜的门放在前面，多数假缝隙在到达昂贵环节之前就被拦下。

其三，选择偏差的账要记全：候选池、方向、参数、以及本文专门提出的一条——观测网格。\textbf{在六个尺度上挑最好的那个，与在六十一条候选里挑最好的一样，都在试验账本上留下一笔。}观测装置的自由度与统计纪律的账本，在这里咬合。

\subsection{三个未完成的方向}

\textbf{分钟级分辨率。}本文窗口检验只覆盖 $\Delta\ge1$ 日一侧；对衰减极快的信号，真实的窗口可能整体落在日内。把 $\Delta_{\min}$ 推到分钟级，需要相应的高频数据、更细的冲击模型与更严格的时点纪律，是这套判据最自然的下一步。

\textbf{容量的内生反馈。}本文把 $\lambda$ 与 $c$ 当作外生参数，但参与规模会同时抬高它们——冲击成本随规模上升，租金耗散随拥挤加速，两个变化同向压缩窗口。把“缝隙的自毁”写成 $\lambda(A,c(A))$ 的动力学模型，是一个具体、有界的理论任务。

\textbf{参照工具库的完备化。}分解公式的检验质量，恰好等于参照工具库的可实施程度。在 A 股约束下（做空受限、工具稀疏），可行的参照系比理论上的贫得多；工具库每补一件可实施的品种，第二道核验就锋利一分。

\subsection{收束}

回到开头那笔小账：信号还在、钱没了。本文想说明的是，这笔账不该只记成本一栏。时间尺度决定你能看见多少，参照系决定这算不算你的，成本决定还剩多少——三栏合起来，才是一个可以用数据结算的问题。\textbf{“缝隙在哪里”的答案从来不是一个地点，而是一组坐标：在什么频率上、相对于什么参照、扣除什么成本之后。}把坐标写全，假缝隙会自己离开，真缝隙才有机会被认真对待。